\documentclass[11pt,a4paper]{article}

\usepackage[utf8]{inputenc}
\usepackage[T1]{fontenc}
\usepackage[spanish]{babel}
\usepackage{amsmath,amssymb}
\usepackage{geometry}
\usepackage{booktabs}
\usepackage{array}
\usepackage{longtable}
\usepackage{hyperref}
\usepackage{xcolor}
\usepackage{enumitem}

\geometry{margin=2.5cm}
\hypersetup{
    colorlinks=true,
    linkcolor=blue,
    urlcolor=blue,
    citecolor=blue
}

\title{DoE para Vibe Coding: Backlog Scrum y Diseño Experimental para un Simulador de Deuda T\'ecnica}
\author{Propuesta de especificaci\'on para proyecto de investigaci\'on}
\date{\today}

\begin{document}

\maketitle

\section{Proyecto 1}

\subsection{Simulador de deuda t\'ecnica extendido con Monte Carlo, optimizaci\'on y persistencia}

El sistema debe permitir simular sprint a sprint la evoluci\'on de:

\begin{equation}
B_k, D_k, V_k, u_k, N_k, R_k, \nu_k
\end{equation}

donde:

\begin{itemize}[leftmargin=*]
    \item $B_k$: backlog funcional remanente.
    \item $D_k$: deuda t\'ecnica remanente.
    \item $V_k$: velocidad efectiva del sprint.
    \item $u_k$: fracci\'on del sprint dedicada a remediaci\'on.
    \item $N_k$: funcionalidad nueva entregada.
    \item $R_k$: deuda remediada.
    \item $\nu_k$: valor econ\'omico generado.
\end{itemize}

\section{\'Epicas del backlog}

\subsection{\'Epica A --- N\'ucleo determin\'istico del modelo}

Implementar el modelo b\'asico de evoluci\'on sprint a sprint.

\subsection{\'Epica B --- Pol\'iticas de decisi\'on}

Implementar distintas estrategias para seleccionar $u_k$.

\subsection{\'Epica C --- Simulaci\'on Monte Carlo}

Ejecutar m\'ultiples corridas variando par\'ametros.

\subsection{\'Epica D --- Optimizaci\'on}

Calcular $u_k$ \textit{\'optimo} bajo restricciones.

\subsection{\'Epica E --- Persistencia y reproducibilidad}

Guardar escenarios, par\'ametros, corridas y resultados.

\subsection{\'Epica F --- Visualizaci\'on y reporting}

Generar CSV, gr\'aficos y reportes comparativos.

\section{Backlog Scrum propuesto}

\subsection{\'Epica A --- N\'ucleo determin\'istico}

\subsubsection{Historia A1 --- Crear estructura de par\'ametros del modelo}

\textbf{Complejidad:} baja.

\textbf{Como} investigador, \textbf{quiero} definir los par\'ametros principales del modelo en una estructura Python, \textbf{para} ejecutar simulaciones reproducibles.

\paragraph{Criterios de aceptaci\'on}

\begin{itemize}[leftmargin=*]
    \item Existe una clase \texttt{ModelParameters}.
    \item Incluye al menos: \texttt{B0}, \texttt{D0}, \texttt{V0}, \texttt{alpha}, \texttt{beta}, \texttt{gamma}, \texttt{theta}, \texttt{lambda\_}, \texttt{rho}, \texttt{K}.
    \item Valida valores negativos no permitidos.
    \item Puede imprimirse o serializarse.
\end{itemize}

\textbf{Riesgo de integraci\'on:} bajo.

\subsubsection{Historia A2 --- Calcular velocidad efectiva $V_k$}

\textbf{Complejidad:} baja.

\textbf{Como} investigador, \textbf{quiero} calcular la velocidad efectiva afectada por deuda, \textbf{para} modelar p\'erdida de productividad.

Ejemplo:

\begin{equation}
V_k = \frac{V_0}{1+\gamma D_k}
\end{equation}

\paragraph{Criterios de aceptaci\'on}

\begin{itemize}[leftmargin=*]
    \item Dado $D_k=0$, se obtiene $V_k=V_0$.
    \item Si $D_k$ aumenta, $V_k$ disminuye.
    \item Hay tests unitarios.
\end{itemize}

\textbf{Riesgo de integraci\'on:} bajo.

\subsubsection{Historia A3 --- Simular un sprint con $u_k$ fijo}

\textbf{Complejidad:} baja-media.

\textbf{Como} investigador, \textbf{quiero} calcular el estado $k+1$ a partir de $B_k,D_k,u_k$, \textbf{para} construir trayectorias de evoluci\'on.

Variables b\'asicas:

\begin{equation}
R_k = u_k V_k
\end{equation}

\begin{equation}
N_k = (1-u_k)V_k
\end{equation}

\begin{equation}
B_{k+1} = \max(0, B_k - N_k)
\end{equation}

\begin{equation}
D_{k+1} = \max(0, D_k - R_k + \alpha N_k + \beta R_k)
\end{equation}

\paragraph{Criterios de aceptaci\'on}

\begin{itemize}[leftmargin=*]
    \item $0 \leq u_k \leq 1$.
    \item $B_k$ nunca queda negativo.
    \item $D_k$ nunca queda negativo.
    \item Se devuelve un objeto \texttt{SprintState}.
\end{itemize}

\textbf{Riesgo de integraci\'on:} medio, porque define el contrato central del simulador.

\subsubsection{Historia A4 --- Simular varios sprints determin\'isticos}

\textbf{Complejidad:} media.

\textbf{Como} investigador, \textbf{quiero} ejecutar una simulaci\'on de $K$ sprints, \textbf{para} observar la evoluci\'on del sistema.

\paragraph{Criterios de aceptaci\'on}

\begin{itemize}[leftmargin=*]
    \item Se genera una tabla con una fila por sprint.
    \item Incluye $B_k,D_k,V_k,u_k,N_k,R_k$.
    \item La simulaci\'on se detiene si $B_k=0$ y $D_k=0$.
    \item Hay tests de regresi\'on simples.
\end{itemize}

\textbf{Riesgo de integraci\'on:} medio.

\subsection{\'Epica B --- Pol\'iticas de decisi\'on}

\subsubsection{Historia B1 --- Pol\'itica naive ``deuda primero''}

\textbf{Complejidad:} baja.

\textbf{Como} investigador, \textbf{quiero} una pol\'itica donde $u_k=1$ mientras exista deuda, \textbf{para} usarla como baseline.

\paragraph{Criterios de aceptaci\'on}

\begin{itemize}[leftmargin=*]
    \item Si $D_k>0$, entonces $u_k=1$.
    \item Si $D_k=0$, entonces $u_k=0$.
    \item Se integra con el simulador de varios sprints.
\end{itemize}

\textbf{Riesgo de integraci\'on:} bajo.

\subsubsection{Historia B2 --- Pol\'itica backlog primero}

\textbf{Complejidad:} baja.

\textbf{Como} investigador, \textbf{quiero} una pol\'itica donde $u_k=0$ mientras exista backlog funcional, \textbf{para} comparar contra deuda primero.

\paragraph{Criterios de aceptaci\'on}

\begin{itemize}[leftmargin=*]
    \item Si $B_k>0$, entonces $u_k=0$.
    \item Si $B_k=0$ y $D_k>0$, entonces $u_k=1$.
    \item Se integra con el simulador.
\end{itemize}

\textbf{Riesgo de integraci\'on:} bajo.

\subsubsection{Historia B3 --- Pol\'itica proporcional}

\textbf{Complejidad:} media.

\textbf{Como} investigador, \textbf{quiero} definir $u_k$ proporcional a la deuda relativa, \textbf{para} tener una pol\'itica heur\'istica intermedia.

Ejemplo:

\begin{equation}
u_k = \frac{D_k}{B_k+D_k}
\end{equation}

\paragraph{Criterios de aceptaci\'on}

\begin{itemize}[leftmargin=*]
    \item $u_k=0$ si $D_k=0$.
    \item $u_k=1$ si $B_k=0$ y $D_k>0$.
    \item $u_k\in[0,1]$.
    \item Maneja correctamente $B_k=D_k=0$.
\end{itemize}

\textbf{Riesgo de integraci\'on:} medio.

\subsubsection{Historia B4 --- Interfaz com\'un de pol\'iticas}

\textbf{Complejidad:} media.

\textbf{Como} desarrollador, \textbf{quiero} que todas las pol\'iticas implementen una interfaz com\'un, \textbf{para} poder intercambiarlas sin modificar el motor de simulaci\'on.

\paragraph{Criterios de aceptaci\'on}

\begin{itemize}[leftmargin=*]
    \item Existe una interfaz \texttt{Policy}.
    \item Cada pol\'itica implementa \texttt{decide\_u(state, params)}.
    \item El motor no depende de clases concretas.
\end{itemize}

\textbf{Riesgo de integraci\'on:} alto, porque refactoriza historias previas.

\subsection{\'Epica C --- Simulaci\'on Monte Carlo}

\subsubsection{Historia C1 --- Definir distribuciones uniformes de par\'ametros}

\textbf{Complejidad:} media.

\textbf{Como} investigador, \textbf{quiero} definir par\'ametros aleatorios con distribuci\'on uniforme, \textbf{para} explorar incertidumbre.

Ejemplo:

\begin{itemize}[leftmargin=*]
    \item $s \sim U(1.0,1.4)$.
    \item $\gamma \sim U(0.00,0.05)$.
    \item $\theta \sim U(0,0.9)$.
    \item $(1-\beta) \sim U(0.5,0.9)$.
    \item $\lambda \sim U(0.2,1.0)$.
\end{itemize}

\paragraph{Criterios de aceptaci\'on}

\begin{itemize}[leftmargin=*]
    \item Se pueden fijar semillas aleatorias.
    \item Se puede muestrear un conjunto completo de par\'ametros.
    \item Hay tests con semilla fija.
\end{itemize}

\textbf{Riesgo de integraci\'on:} medio.

\subsubsection{Historia C2 --- Ejecutar $N$ simulaciones Monte Carlo}

\textbf{Complejidad:} media.

\textbf{Como} investigador, \textbf{quiero} ejecutar muchas corridas, \textbf{para} estimar distribuciones de resultados.

\paragraph{Criterios de aceptaci\'on}

\begin{itemize}[leftmargin=*]
    \item Permite configurar \texttt{n\_runs}.
    \item Devuelve resultados agregados.
    \item Conserva los resultados individuales.
    \item Usa una pol\'itica seleccionable.
\end{itemize}

\textbf{Riesgo de integraci\'on:} alto, porque combina modelo, pol\'iticas y par\'ametros aleatorios.

\subsubsection{Historia C3 --- Agregar m\'etricas de salida}

\textbf{Complejidad:} media.

\textbf{Como} investigador, \textbf{quiero} calcular m\'etricas por corrida, \textbf{para} comparar estrategias.

M\'etricas sugeridas:

\begin{itemize}[leftmargin=*]
    \item Sprints hasta finalizar backlog.
    \item Sprints hasta eliminar deuda.
    \item Deuda final.
    \item Valor acumulado descontado.
    \item $u_k$ promedio.
    \item Varianza de $u_k$.
    \item Productividad efectiva promedio.
    \item Defectos/inconsistencias detectadas.
\end{itemize}

\paragraph{Criterios de aceptaci\'on}

\begin{itemize}[leftmargin=*]
    \item Las m\'etricas se calculan para cada corrida.
    \item Se puede obtener media, desv\'io est\'andar y percentiles.
    \item Se exportan a CSV.
\end{itemize}

\textbf{Riesgo de integraci\'on:} alto.

\subsection{\'Epica D --- Optimizaci\'on}

\subsubsection{Historia D1 --- Implementar b\'usqueda por grilla de $u_k$}

\textbf{Complejidad:} media.

\textbf{Como} investigador, \textbf{quiero} evaluar muchos valores posibles de $u_k$, \textbf{para} encontrar una decisi\'on aproximadamente \textit{\'optima}.

Ejemplo:

\begin{equation}
u_k \in \{0.0,0.01,0.02,\dots,1.0\}
\end{equation}

\paragraph{Criterios de aceptaci\'on}

\begin{itemize}[leftmargin=*]
    \item Eval\'ua una funci\'on objetivo.
    \item Respeta $0 \leq u_k \leq 1$.
    \item Devuelve el mejor $u_k$.
    \item Maneja el caso $D_k=0 \Rightarrow u_k=0$.
\end{itemize}

\textbf{Riesgo de integraci\'on:} medio-alto.

\subsubsection{Historia D2 --- Definir funci\'on objetivo econ\'omica}

\textbf{Complejidad:} alta.

\textbf{Como} investigador, \textbf{quiero} optimizar $u_k$ seg\'un una funci\'on econ\'omica, \textbf{para} balancear entrega de valor y reducci\'on de deuda.

Ejemplo conceptual:

\begin{equation}
J(u_k)=
\text{valor entregado}
+
\text{valor de capacidad futura}
-
\text{costo de demora}
-
\text{costo de deuda residual}
\end{equation}

\paragraph{Criterios de aceptaci\'on}

\begin{itemize}[leftmargin=*]
    \item La funci\'on objetivo es configurable.
    \item Usa $B_k,D_k,V_k,\lambda,\rho,\theta$.
    \item Puede compararse contra pol\'iticas naive.
    \item Hay tests de casos l\'imite.
\end{itemize}

\textbf{Riesgo de integraci\'on:} alto.

\subsubsection{Historia D3 --- Pol\'itica \textit{\'optima} local}

\textbf{Complejidad:} alta.

\textbf{Como} investigador, \textbf{quiero} una pol\'itica que calcule $u_k$ \textit{\'optimo} en cada sprint, \textbf{para} comparar decisiones \textit{\'optimas} contra heur\'isticas.

\paragraph{Criterios de aceptaci\'on}

\begin{itemize}[leftmargin=*]
    \item Implementa la interfaz \texttt{Policy}.
    \item Calcula $u_k$ usando b\'usqueda por grilla.
    \item Respeta restricciones de frontera.
    \item Se integra con Monte Carlo.
\end{itemize}

\textbf{Riesgo de integraci\'on:} muy alto.

\subsubsection{Historia D4 --- Comparaci\'on entre pol\'iticas}

\textbf{Complejidad:} alta.

\textbf{Como} investigador, \textbf{quiero} ejecutar el mismo escenario con varias pol\'iticas, \textbf{para} comparar resultados.

\paragraph{Criterios de aceptaci\'on}

\begin{itemize}[leftmargin=*]
    \item Ejecuta naive, backlog-first, proporcional y \textit{\'optima}.
    \item Produce m\'etricas comparables.
    \item Exporta una tabla por pol\'itica.
    \item Genera gr\'aficos comparativos.
\end{itemize}

\textbf{Riesgo de integraci\'on:} muy alto.

\subsection{\'Epica E --- Persistencia}

\subsubsection{Historia E1 --- Guardar resultados en CSV}

\textbf{Complejidad:} baja.

\textbf{Como} investigador, \textbf{quiero} guardar resultados en CSV, \textbf{para} analizarlos externamente.

\paragraph{Criterios de aceptaci\'on}

\begin{itemize}[leftmargin=*]
    \item Exporta estados sprint a sprint.
    \item Exporta m\'etricas finales.
    \item Los nombres de columnas son estables.
\end{itemize}

\textbf{Riesgo de integraci\'on:} bajo.

\subsubsection{Historia E2 --- Guardar escenarios en JSON}

\textbf{Complejidad:} media.

\textbf{Como} investigador, \textbf{quiero} guardar par\'ametros y configuraci\'on experimental en JSON, \textbf{para} repetir una corrida.

\paragraph{Criterios de aceptaci\'on}

\begin{itemize}[leftmargin=*]
    \item Guarda par\'ametros.
    \item Guarda pol\'itica usada.
    \item Guarda semilla aleatoria.
    \item Puede recargarse el escenario.
\end{itemize}

\textbf{Riesgo de integraci\'on:} medio.

\subsubsection{Historia E3 --- Reproducibilidad completa}

\textbf{Complejidad:} alta.

\textbf{Como} investigador, \textbf{quiero} reproducir una corrida completa desde archivo, \textbf{para} validar resultados experimentales.

\paragraph{Criterios de aceptaci\'on}

\begin{itemize}[leftmargin=*]
    \item Un archivo JSON permite relanzar una simulaci\'on.
    \item Con la misma semilla produce los mismos resultados.
    \item Registra versi\'on del modelo.
    \item Registra fecha/hora de ejecuci\'on.
\end{itemize}

\textbf{Riesgo de integraci\'on:} alto.

\subsection{\'Epica F --- Visualizaci\'on}

\subsubsection{Historia F1 --- Graficar evoluci\'on de $B_k$ y $D_k$}

\textbf{Complejidad:} baja.

\textbf{Como} investigador, \textbf{quiero} visualizar backlog y deuda por sprint, \textbf{para} interpretar trayectorias.

\paragraph{Criterios de aceptaci\'on}

\begin{itemize}[leftmargin=*]
    \item Genera PNG.
    \item Muestra $B_k$ y $D_k$.
    \item Permite elegir corrida o promedio.
\end{itemize}

\textbf{Riesgo de integraci\'on:} bajo.

\subsubsection{Historia F2 --- Boxplot de $u_k$ \textit{\'optimo}}

\textbf{Complejidad:} media.

\textbf{Como} investigador, \textbf{quiero} visualizar la distribuci\'on de $u_k$, \textbf{para} observar variabilidad bajo incertidumbre.

\paragraph{Criterios de aceptaci\'on}

\begin{itemize}[leftmargin=*]
    \item Genera boxplot.
    \item Agrupa por escenario.
    \item Exporta PNG.
\end{itemize}

\textbf{Riesgo de integraci\'on:} medio.

\subsubsection{Historia F3 --- Heatmap de sensibilidad}

\textbf{Complejidad:} alta.

\textbf{Como} investigador, \textbf{quiero} generar heatmaps variando dos par\'ametros, \textbf{para} observar regiones de decisi\'on.

Ejemplo:

\begin{itemize}[leftmargin=*]
    \item Eje X: $B_k/(M s)$.
    \item Eje Y: $D_k/B_k$.
    \item Color: $u_k$ \textit{\'optimo} promedio.
\end{itemize}

\paragraph{Criterios de aceptaci\'on}

\begin{itemize}[leftmargin=*]
    \item Permite elegir dos par\'ametros de barrido.
    \item Calcula promedio de $u_k$.
    \item Exporta imagen y CSV.
    \item Integra Monte Carlo + optimizaci\'on.
\end{itemize}

\textbf{Riesgo de integraci\'on:} muy alto.

\section{Backlog resumido por complejidad}

\subsection{Historias bajas}

\begin{table}[h!]
\centering
\begin{tabular}{ll}
\toprule
\textbf{ID} & \textbf{Historia} \\
\midrule
A1 & Par\'ametros del modelo \\
A2 & Velocidad efectiva \\
B1 & Pol\'itica deuda primero \\
B2 & Pol\'itica backlog primero \\
E1 & Exportar CSV \\
F1 & Gr\'afico $B_k,D_k$ \\
\bottomrule
\end{tabular}
\end{table}

\subsection{Historias medias}

\begin{table}[h!]
\centering
\begin{tabular}{ll}
\toprule
\textbf{ID} & \textbf{Historia} \\
\midrule
A3 & Simular un sprint \\
A4 & Simular varios sprints \\
B3 & Pol\'itica proporcional \\
B4 & Interfaz com\'un de pol\'iticas \\
C1 & Distribuciones uniformes \\
C2 & Corridas Monte Carlo \\
C3 & M\'etricas \\
D1 & B\'usqueda por grilla \\
E2 & Escenarios JSON \\
F2 & Boxplots \\
\bottomrule
\end{tabular}
\end{table}

\subsection{Historias altas}

\begin{table}[h!]
\centering
\begin{tabular}{ll}
\toprule
\textbf{ID} & \textbf{Historia} \\
\midrule
D2 & Funci\'on objetivo econ\'omica \\
D3 & Pol\'itica \textit{\'optima} local \\
D4 & Comparaci\'on de pol\'iticas \\
E3 & Reproducibilidad completa \\
F3 & Heatmap de sensibilidad \\
\bottomrule
\end{tabular}
\end{table}

\section{Orden sugerido de integraci\'on}

Para generar problemas de interacci\'on medibles, no conviene integrar todo linealmente. Conviene usar secuencias alternativas.

\subsection{Secuencia base}

\begin{enumerate}[leftmargin=*]
    \item A1 --- Par\'ametros.
    \item A2 --- Velocidad.
    \item A3 --- Un sprint.
    \item A4 --- Simulaci\'on determin\'istica.
    \item B1 --- Deuda primero.
    \item B2 --- Backlog primero.
    \item B3 --- Proporcional.
    \item B4 --- Interfaz com\'un.
    \item C1 --- Distribuciones.
    \item C2 --- Monte Carlo.
    \item C3 --- M\'etricas.
    \item D1 --- Grilla.
    \item D2 --- Funci\'on objetivo.
    \item D3 --- Pol\'itica \textit{\'optima}.
    \item D4 --- Comparaci\'on.
    \item E1/E2/E3 --- Persistencia.
    \item F1/F2/F3 --- Visualizaci\'on.
\end{enumerate}

\subsection{Secuencia alternativa para inducir fricci\'on}

\begin{enumerate}[leftmargin=*]
    \item A1.
    \item A2.
    \item B1.
    \item B2.
    \item E1.
    \item A3.
    \item A4.
    \item C1.
    \item F1.
    \item B4.
    \item C2.
    \item C3.
    \item D1.
    \item D2.
    \item D3.
    \item F2.
    \item F3.
    \item E3.
\end{enumerate}

Esta segunda secuencia fuerza refactorizaciones porque introduce pol\'iticas y exportaci\'on antes de consolidar el motor.

\section{Dise\~no de Experimentos factorial}

\subsection{Objetivo del DoE}

Medir c\'omo distintas condiciones de desarrollo afectan:

\begin{itemize}[leftmargin=*]
    \item Productividad.
    \item Defectos funcionales.
    \item Defectos de integraci\'on.
    \item Retrabajo.
    \item Estabilidad del c\'odigo.
    \item Dificultad de integraci\'on entre historias.
    \item Calidad del resultado generado por Vibe Coding.
\end{itemize}

\section{Factores experimentales}

\subsection{Factor A --- Complejidad de historia}

\begin{table}[h!]
\centering
\begin{tabular}{ll}
\toprule
\textbf{Nivel} & \textbf{Descripci\'on} \\
\midrule
A1 & Baja \\
A2 & Media \\
A3 & Alta \\
\bottomrule
\end{tabular}
\end{table}

\subsection{Factor B --- Tipo de acoplamiento dominante}

\begin{table}[h!]
\centering
\begin{tabular}{ll}
\toprule
\textbf{Nivel} & \textbf{Descripci\'on} \\
\midrule
B1 & Acoplamiento de datos \\
B2 & Acoplamiento de control \\
B3 & Acoplamiento temporal \\
B4 & Acoplamiento matem\'atico \\
\bottomrule
\end{tabular}
\end{table}

Ejemplos:

\begin{itemize}[leftmargin=*]
    \item Datos: CSV, JSON, estructuras compartidas.
    \item Control: pol\'iticas que modifican el flujo.
    \item Temporal: Monte Carlo, semillas, orden de ejecuci\'on.
    \item Matem\'atico: funci\'on objetivo, optimizaci\'on, restricciones.
\end{itemize}

\subsection{Factor C --- Orden de integraci\'on}

\begin{table}[h!]
\centering
\begin{tabular}{ll}
\toprule
\textbf{Nivel} & \textbf{Descripci\'on} \\
\midrule
C1 & Orden incremental natural \\
C2 & Orden con refactorizaci\'on forzada \\
C3 & Orden aleatorizado controlado \\
\bottomrule
\end{tabular}
\end{table}

\subsection{Factor D --- Modalidad de desarrollo}

\begin{table}[h!]
\centering
\begin{tabular}{ll}
\toprule
\textbf{Nivel} & \textbf{Descripci\'on} \\
\midrule
D1 & Programaci\'on manual tradicional \\
D2 & Vibe Coding con prompts simples \\
D3 & Vibe Coding con prompts estructurados \\
D4 & Vibe Coding con tests previos \\
\bottomrule
\end{tabular}
\end{table}

\subsection{Factor E --- Nivel de especificaci\'on}

\begin{table}[h!]
\centering
\begin{tabular}{ll}
\toprule
\textbf{Nivel} & \textbf{Descripci\'on} \\
\midrule
E1 & Historia breve \\
E2 & Historia + criterios de aceptaci\'on \\
E3 & Historia + criterios + tests esperados \\
E4 & Historia + contrato de interfaces \\
\bottomrule
\end{tabular}
\end{table}

\section{Dise\~no factorial m\'inimo recomendado}

Un factorial completo ser\'ia demasiado grande:

\begin{equation}
3 \times 4 \times 3 \times 4 \times 4 = 576
\end{equation}

Eso es excesivo. Se recomienda un dise\~no fraccional o por bloques.

\subsection{Dise\~no inicial recomendado}

Usar estos factores principales:

\begin{table}[h!]
\centering
\begin{tabular}{ll}
\toprule
\textbf{Factor} & \textbf{Niveles} \\
\midrule
Complejidad & baja, media, alta \\
Modalidad & manual, vibe simple, vibe estructurado, vibe con tests \\
Nivel de especificaci\'on & breve, criterios, tests, contrato \\
\bottomrule
\end{tabular}
\end{table}

Total:

\begin{equation}
3 \times 4 \times 4 = 48
\end{equation}

Con 2 repeticiones:

\begin{equation}
48 \times 2 = 96 \text{ ejecuciones}
\end{equation}

\section{Variables dependientes}

\subsection{Productividad}

\begin{table}[h!]
\centering
\begin{tabular}{ll}
\toprule
\textbf{M\'etrica} & \textbf{Descripci\'on} \\
\midrule
Tiempo de implementaci\'on & Minutos por historia \\
Tiempo hasta primer c\'odigo ejecutable & Minutos \\
Tiempo hasta pasar tests & Minutos \\
Cantidad de iteraciones de prompt & N\'umero \\
\bottomrule
\end{tabular}
\end{table}

\subsection{Calidad}

\begin{table}[h!]
\centering
\begin{tabular}{ll}
\toprule
\textbf{M\'etrica} & \textbf{Descripci\'on} \\
\midrule
Defectos unitarios & Errores dentro de una historia \\
Defectos de integraci\'on & Errores al combinar historias \\
Defectos matem\'aticos & Inconsistencias del modelo \\
Defectos de borde & Errores con $B_k=0$, $D_k=0$, $u_k=0/1$ \\
\bottomrule
\end{tabular}
\end{table}

\subsection{Mantenibilidad}

\begin{table}[h!]
\centering
\begin{tabular}{ll}
\toprule
\textbf{M\'etrica} & \textbf{Descripci\'on} \\
\midrule
Complejidad ciclom\'atica & Complejidad estructural del c\'odigo \\
Duplicaci\'on & Repetici\'on de fragmentos o estructuras \\
Tama\~no de archivos & Longitud y distribuci\'on del c\'odigo \\
Cantidad de refactorizaciones & Cambios estructurales necesarios \\
Cambios necesarios para integrar nueva historia & Medida de esfuerzo de integraci\'on \\
\bottomrule
\end{tabular}
\end{table}

\subsection{Robustez}

\begin{table}[h!]
\centering
\begin{tabular}{ll}
\toprule
\textbf{M\'etrica} & \textbf{Descripci\'on} \\
\midrule
Tests pasados & Proporci\'on o cantidad de pruebas exitosas \\
Cobertura & Cobertura de tests \\
Fallos ante valores extremos & Robustez frente a condiciones de borde \\
Reproducibilidad con semilla fija & Capacidad de repetir resultados \\
\bottomrule
\end{tabular}
\end{table}

\section{Hip\'otesis experimentales}

\subsection{H1 --- Productividad}

El Vibe Coding con prompts estructurados reduce el tiempo de implementaci\'on frente a programaci\'on manual.

\subsection{H2 --- Calidad}

El Vibe Coding con tests previos reduce defectos de integraci\'on frente a Vibe Coding con prompts simples.

\subsection{H3 --- Complejidad}

La ventaja de productividad del Vibe Coding disminuye en historias de alta complejidad matem\'atica.

\subsection{H4 --- Especificaci\'on}

Las historias con contrato de interfaces producen menos defectos de integraci\'on que las historias breves.

\subsection{H5 --- Interacci\'on}

Existe interacci\'on entre modalidad de desarrollo y nivel de especificaci\'on:

\begin{equation}
\text{Vibe Coding simple} + \text{historia breve}
\end{equation}

producir\'a m\'as errores que:

\begin{equation}
\text{Vibe Coding estructurado} + \text{contrato de interfaces}
\end{equation}

\section{Modelo estad\'istico sugerido}

Para analizar resultados:

\begin{equation}
Y =
\mu
+ A_i
+ D_j
+ E_k
+ (A D)_{ij}
+ (A E)_{ik}
+ (D E)_{jk}
+ \varepsilon
\end{equation}

donde:

\begin{itemize}[leftmargin=*]
    \item $Y$: m\'etrica observada.
    \item $A_i$: efecto de complejidad.
    \item $D_j$: efecto de modalidad de desarrollo.
    \item $E_k$: efecto de nivel de especificaci\'on.
    \item $(AD)$, $(AE)$, $(DE)$: interacciones.
    \item $\varepsilon$: error experimental.
\end{itemize}

M\'etricas posibles para $Y$:

\begin{itemize}[leftmargin=*]
    \item Tiempo de implementaci\'on.
    \item N\'umero de defectos.
    \item Defectos de integraci\'on.
    \item N\'umero de prompts.
    \item Cobertura alcanzada.
    \item Esfuerzo de retrabajo.
\end{itemize}

\section{Unidad experimental}

La unidad experimental recomendada es:

\begin{quote}
Una historia implementada bajo una combinaci\'on concreta de complejidad, modalidad de desarrollo y nivel de especificaci\'on.
\end{quote}

Ejemplo:

\begin{table}[h!]
\centering
\begin{tabular}{llll}
\toprule
\textbf{Historia} & \textbf{Complejidad} & \textbf{Modalidad} & \textbf{Especificaci\'on} \\
\midrule
D3 & Alta & Vibe con tests & Contrato de interfaces \\
\bottomrule
\end{tabular}
\end{table}

\section{Resultado esperado del experimento}

El estudio deber\'ia permitir responder:

\begin{enumerate}[leftmargin=*]
    \item \textquestiondown Vibe Coding acelera la implementaci\'on?
    \item \textquestiondown D\'onde introduce defectos?
    \item \textquestiondown Qu\'e historias son m\'as riesgosas?
    \item \textquestiondown Qu\'e nivel de especificaci\'on reduce errores?
    \item \textquestiondown Cu\'ando aparecen problemas de integraci\'on?
    \item \textquestiondown La generaci\'on asistida funciona mejor con tests previos?
    \item \textquestiondown Qu\'e tipo de acoplamiento degrada m\'as la productividad?
\end{enumerate}

\section{Recomendaci\'on final}

Para la investigaci\'on, este backlog debe usarse como instrumento experimental y no solo como especificaci\'on funcional.

La combinaci\'on m\'as valiosa ser\'ia:

\begin{equation}
\text{Complejidad}
\times
\text{Modalidad de desarrollo}
\times
\text{Nivel de especificaci\'on}
\end{equation}

con foco especial en medir:

\begin{equation}
\text{defectos de integraci\'on}
\end{equation}

porque ah\'i es donde probablemente aparezca la se\~nal m\'as interesante del experimento.

\end{document}
